\subsection{\texorpdfstring{第~\ref{sec:isometric-immersions}~节}{等距浸入}}\label{subsec:solutions-isometric-immersions}

\begin{description}[font=\normalfont,leftmargin=2em,labelwidth=1.5em]
\RGsolution{ex:second-fundamental-form-properties}{second-fundamental-form-properties}
\begin{enumerate}[label=(\alph*),ref=\alph*,leftmargin=*]
\item 这由如下事实可见: $(\widetilde\nabla_{\widetilde X}\widetilde Y)|_N$ 只依赖于 $\widetilde X|_N=X$ 和 $\widetilde Y|_N=Y$.
\item 这也直接来自
\[
  B(X,Y)=\widetilde\nabla_{\widetilde X}\widetilde Y
           -(\widetilde\nabla_{\widetilde X}\widetilde Y)^\top
         =(\widetilde\nabla_{\widetilde X}\widetilde Y)^\perp.
\]
\item 由第~\ref{chap:differentiable-manifolds}~章的 Exercise~\ref{ex:vector-fields}(\ref{ex:lie-bracket-tangent-to-submanifold}) 知道, $[\widetilde X,\widetilde Y]$ 与 $N$ 相切. 因此,
\[
  B(X,Y)-B(Y,X)=(\widetilde\nabla_{\widetilde X}\widetilde Y
                       -\widetilde\nabla_{\widetilde Y}\widetilde X)^\perp
              =[\widetilde X,\widetilde Y]^\perp=0.
\]
\item 设 $f$ 是 $N$ 上的函数, $\widetilde f$ 是 $f$ 到 $M$ 上的任一延拓. 则
\[
  B(fX,Y)=(\widetilde\nabla_{\widetilde f\widetilde X}\widetilde Y)^\perp
         =f(\widetilde\nabla_{\widetilde X}\widetilde Y)^\perp=fB(X,Y),
\]
且由对称性,
\[
  B(X,fY)=fB(X,Y).
\]
此外, 显然有
\[
  B(X+Y,Z)=B(X,Z)+B(Y,Z)
\]
和
\[
  B(X,Y+Z)=B(X,Y)+B(X,Z).
\]
\item 若 $x:V\to\mathbb{R}^n$ 是 $p\in M$ 附近的局部坐标, 则
\[
  B(X,Y)=B\left(\sum_{i=1}^nX^i\frac{\partial}{\partial x^i},
               \sum_{j=1}^nY^j\frac{\partial}{\partial x^j}\right)
        =\sum_{i,j=1}^nX^iY^jB\left(\frac{\partial}{\partial x^i},\frac{\partial}{\partial x^j}\right).
\]
因此,
\[
  B(X,Y)_p=\sum_{i,j=1}^nX^i(p)Y^j(p)B\left(\frac{\partial}{\partial x^i},
                                           \frac{\partial}{\partial x^j}\right)_p
\]
只依赖于 $X_p,Y_p$.
\item 由第~\ref{chap:differentiable-manifolds}~章的 Exercise~\ref{ex:vector-fields}(\ref{ex:lie-bracket-tangent-to-submanifold}) 可见, $[\widetilde X,\widetilde Y]$ 是 $[X,Y]$ 的延拓. 因而
\[
  \widetilde\nabla_{[\widetilde X,\widetilde Y]}\widetilde X-\nabla_{[X,Y]}X=B([X,Y],X)
\]
与 $N$ 正交.
\end{enumerate}

\RGsolution{ex:gauss-curvature-as-ratio-of-volumes}{gauss-curvature-as-ratio-of-volumes}
设 $g:D\subset N\to S^n$ 是定义在 $p$ 的邻域 $D$ 上的 Gauss 映射. 由于 $\det(\mathrm{d}g)_p=(-1)^n\det(S_{n_p})=(-1)^nK(p)\ne0$, 可以假定 $g$ 在 $D$ 上是微分同胚. 若 $\omega$ 是 $S^n$ 的标准体积形式, 则
\[
  \operatorname{vol}(g(D))=\int_{g(D)}\omega=\left|\int_D g^*\omega\right|
\]
(因为 $g$ 可能反转定向). 若 $\{v_1,\ldots,v_n\}$ 是 $T_qN\cong T_{g(q)}S^n$ 的一组标准正交基, 且 $q\in D$, 则
\begin{align*}
  g^*\omega(v_1,\ldots,v_n)
  &=\omega((\mathrm{d}g)_qv_1,\ldots,(\mathrm{d}g)_qv_n)\\
  &=\det(\mathrm{d}g)_q\,\omega(v_1,\ldots,v_n)=\det(\mathrm{d}g)_q,
\end{align*}
所以在 $D$ 上 $g^*\omega=\det(\mathrm{d}g)\omega_N$ (其中 $\omega_N$ 是 $N$ 的体积元). 因此
\[
  \operatorname{vol}(g(D))=\left|\int_D\det(\mathrm{d}g)\right|
    =\int_D|\det(\mathrm{d}g)|=\int_D|K|
\]
(因为 $\det(\mathrm{d}g)$ 在 $D$ 上不变号). 由中值定理即得所求结果:
\[
  |K(p)|=\lim_{D\to p}\frac{\int_D|K|}{\operatorname{vol}(D)}
         =\lim_{D\to p}\frac{\operatorname{vol}(g(D))}{\operatorname{vol}(D)}.
\]

\RGsolution{ex:sectional-curvature-by-geodesic-surface}{sectional-curvature-by-geodesic-surface}
\begin{enumerate}[label=(\alph*),ref=\alph*,leftmargin=*]
\item 球 $B_\varepsilon(p)$ 由从 $p$ 出发且长度为 $\varepsilon$ 的测地线覆盖, 这些测地线由 $c_v(t):=\exp_p(tv)$ 给出, 其中 $0\le t\le\varepsilon$ 且 $\lVert v\rVert=1$. 因此 $N_p:=\exp_p(B_\varepsilon\cap\Pi)$ 由 $B_\varepsilon(p)$ 中与 $\Pi$ 相切的测地线构成. 选取 $T_pM$ 的标准正交基 $\{(E_1)_p,\ldots,(E_n)_p\}$, 使得 $\{(E_1)_p,(E_2)_p\}$ 是 $\Pi$ 的一组基. 取正规坐标
\[
  \varphi(x^1,\ldots,x^n)=\exp_p\left(\sum_{i=1}^n x^i(E_i)_p\right),
\]
有
\[
  N_p=\{q\in B_\varepsilon(p)\mid x^3(q)=\cdots=x^n(q)=0\}.
\]
由第~\ref{chap:differentiable-manifolds}~章的 Exercise~\ref{ex:immersions-embeddings}(\ref{ex:coordinate-submanifold-characterization}) 可知 $N_p$ 是 $M$ 的 2 维子流形.
\item 由第~\ref{chap:riemannian-manifolds}~章的 Exercise~\ref{ex:minimizing-properties-of-geodesics}(\ref{ex:metric-and-connection-in-normal-coordinates}) 知道, 对 $i,j=1,\ldots,n$, 正规坐标满足
\[
  \left(\widetilde\nabla_{\frac{\partial}{\partial x^i}}\frac{\partial}{\partial x^j}\right)_p=0,
\]
其中 $\widetilde\nabla$ 是 $M$ 的 Levi–Civita 联络. 因此, 对 $i,j=1,2$, 有
\[
  B\left(\left(\frac{\partial}{\partial x^i}\right)_p,
          \left(\frac{\partial}{\partial x^j}\right)_p\right)
     =\left(\widetilde\nabla_{\frac{\partial}{\partial x^i}}
             \frac{\partial}{\partial x^j}\right)_p^\perp=0,
\]
所以 $N_p$ 的第二基本形式在 $p$ 处为零. 从 Proposition~\ref{prop:sectional-curvature-of-isometrically-immersed-submanifold} 得
\[
  K^{N_p}(\Pi)=K^M(\Pi).
\]
\end{enumerate}

\RGsolution{ex:principal-curvatures-from-geodesic-curvatures}{principal-curvatures-from-geodesic-curvatures}
设 $c:I\to N$ 是 $N$ 中以弧长为参数的测地线, 它在 $c(s)$ 处与主方向 $(E_i)_{c(s)}$ (即 $S_{n_{c(s)}}$ 的一个单位特征向量) 相切. 则 $c$ (在 $M$ 中) 的测地曲率由下式给出:
\begin{align*}
  k_g(s)&=\lVert\widetilde\nabla_{\dot c(s)}\dot c(s)\rVert
     =\lVert B(\dot c(s),\dot c(s))-\nabla_{\dot c(s)}\dot c(s)\rVert\\
  &=\lVert B(\dot c(s),\dot c(s))\rVert
    =\bigl|\langle S_{n_{c(s)}}(\dot c(s)),\dot c(s)\rangle\bigr|\\
  &=\bigl|\langle S_{n_{c(s)}}((E_i)_{c(s)}),(E_i)_{c(s)}\rangle\bigr|
    =|\lambda_i|.
\end{align*}

\RGsolution{ex:gauss-maps-of-paraboloid-and-saddle}{gauss-maps-of-paraboloid-and-saddle}
\begin{enumerate}[label=(\alph*),ref=\alph*,leftmargin=*]
\item 抛物面的一个参数化例如由下式给出的映射 $\varphi:\mathbb{R}^2\to\mathbb{R}^3$:
\[
  \varphi(v,v)=\left(v,v,\frac12(v^2+v^2)\right).
\]
抛物面在 $\varphi(v,v)$ 处的切空间由
\[
  \frac{\partial\varphi}{\partial v}=(1,0,v)
   \quad\text{和}\quad\frac{\partial\varphi}{\partial v}=(0,1,v)
\]
生成, 一个法向量是
\[
  \frac{\partial\varphi}{\partial v}\times\frac{\partial\varphi}{\partial v}=(-v,-v,1).
\]
因此, Gauss 映射的一种选取为
\[
  g(v,v)=\frac{(v,v,-1)}{\sqrt{1+v^2+v^2}}.
\]
注意抛物面在这个映射下的像包含在 $S^2$ 的南半球, 在那里仍可用 $\mathbb{R}^3$ 的前两个坐标函数作为坐标. 在这些坐标下, Gauss 映射可直接写为
\[
  \widehat g(v,v)=\frac{(v,v)}{\sqrt{1+v^2+v^2}},
\]
其导数由 Jacobian 矩阵
\[
  \mathrm{d}\widehat g=\begin{pmatrix}
  (1+v^2)(1+v^2+v^2)^{-3/2}&-vv(1+v^2+v^2)^{-3/2}\\
  -vv(1+v^2+v^2)^{-3/2}&(1+v^2)(1+v^2+v^2)^{-3/2}
  \end{pmatrix}
\]
表示. 因而抛物面的 Gauss 曲率为
\[
  K=\det(\mathrm{d}g)=(1+v^2+v^2)^{-2}.
\]
\item 鞍面的一个参数化例如由下式给出的映射 $\varphi:\mathbb{R}^2\to\mathbb{R}^2$:
\[
  \varphi(v,v)=(v,v,vv).
\]
鞍面在 $\varphi(v,v)$ 处的切空间由
\[
  \frac{\partial\varphi}{\partial v}=(1,0,v)
   \quad\text{和}\quad\frac{\partial\varphi}{\partial v}=(0,1,v)
\]
生成, 一个法向量是
\[
  \frac{\partial\varphi}{\partial v}\times\frac{\partial\varphi}{\partial v}=(-v,-v,1).
\]
因此, Gauss 映射的一种选取为
\[
  g(v,v)=\frac{(v,v,-1)}{\sqrt{1+v^2+v^2}}.
\]
注意鞍面在这个映射下的像包含在 $S^2$ 的南半球, 在那里仍可用 $\mathbb{R}^3$ 的前两个坐标函数作为坐标. 在这些坐标下, Gauss 映射可直接写为
\[
  \widehat g(v,v)=\frac{(v,v)}{\sqrt{1+v^2+v^2}},
\]
其导数由 Jacobian 矩阵
\[
  \mathrm{d}\widehat g=\begin{pmatrix}
  -vv(1+v^2+v^2)^{-3/2}&(1+v^2)(1+v^2+v^2)^{-3/2}\\
  (1+v^2)(1+v^2+v^2)^{-3/2}&-vv(1+v^2+v^2)^{-3/2}
  \end{pmatrix}
\]
表示. 因而鞍面的 Gauss 曲率为
\[
  K=\det(\mathrm{d}g)=-(1+v^2+v^2)^{-2}.
\]
\end{enumerate}
\end{description}
